\section{Conventions, definitions and the frozen statements}\label{bpN:p0}

This part fixes notation and restates, in mathematical form, the three frozen theorems of track N
(owner approval D20, 27 October 2026) and the eleven frozen definitions of
\texttt{\seqsplit{KAM/Frozen/Nekhoroshev/Defs/}} together with the five track-K definitions they use. The Lean
declarations are authoritative; the statements below are transcriptions of them, with the Lean names
recorded under each. N-HEAD with real bounds only is \emph{not} frozen (source gap F-N1, D17/D18/D20); the
frozen headline is N-HEAD-C. The additions of D22 (N-PARTIAL-CONVEX, N-AVERAGING) were frozen by D25 and D28 and are transcribed below. Results of track K's blueprint are cited as \Kref{bp:p1:\dots}.

\subsection{Conventions}\label{bpN:p0:conventions}

\begin{itemize}
\item $n\in\N$ is the number of degrees of freedom; N-HEAD-C assumes $n\ge3$ (D17), N-ZERO and N-QC are stated
  for every $n$. A phase point is $z=(\theta,I)$ with angles $\theta$ and actions $I$, in this order.
\item Ambient spaces are those of the Lean text: $\R^n=(\texttt{Fin }n\to\R)$, $\C^n=(\texttt{Fin }n\to\C)$,
  with the max norms $\abs v=\max_i\abs{v_i}$ on both. The embedding $\R^n\subset\C^n$ is componentwise. The
  norm $\norm{\cdot}$ in every frozen statement is the max norm; on the phase space $\R^n\times\R^n$ it is
  $\norm{(\theta,I)}=\max(\abs\theta,\abs I)$ (Lean product norm). The only Euclidean structure is the dot product
  $\sum_i\omega_iv_i$ used for orthogonality in quasi-convexity; Euclidean norms used later for comparison with the
  sources come with explicit inequalities ($\abs v\le\norm v_2\le\sqrt n\abs v$). For $n=0$ all these spaces are
  $\{0\}$ and the norm is $0$.
\item Operator norms of linear and multilinear maps are the induced norms of these max norms. In particular
  $\norm{D^2h(z)}$ for complex $h$ is the norm of $D^2h(z)\in\mathcal L(\C^n,\mathcal L(\C^n,\C))$ (Lean
  \Lean{fderiv} applied twice), i.e.\ $\sup_{\abs u,\abs v\le1}\abs{D^2h(z)[u,v]}$.
\item Angles are \emph{not} reduced modulo $\Z^n$ in the statements: $\theta\in\R^n$ or $\theta\in T_s\subset\C^n$,
  and $1$-periodicity is a hypothesis on $f$ (Definition~\ref{bpN:p0:pert}). So the torus is $\R^n/\Z^n$ and
  the Fourier characters are $e^{2\pi ik\cdot\theta}$ as in track K (\Kref{bp:p1:decay},
  \Kref{bp:p1:summation}). P\"oschel's torus is $\R^n/2\pi\Z^n$ and his norms are Euclidean; the conversion is
  a separate step of the proof (Part~1 of \texttt{PLAN.md}).
\item The complex strip is $T_s=\{\theta\in\C^n:\abs{\operatorname{Im}\theta_i}<s\ \forall i\}$ (Lean
  \Lean{KAM.complexStrip}), as in track K.
\end{itemize}

\subsection{Track-K definitions used}\label{bpN:p0:kdefs}

\begin{definition}[Complex action neighbourhood]\label{bpN:p0:actionnbhd}
For $B\subset\R^n$ and $r\in\R$, $B_r=\{I\in\C^n:\exists J\in B,\ \abs{I-J}<r\}$ (max norm, $J$ viewed in $\C^n$).
Thus $B_r=\emptyset$ for $r\le0$, and for $B$ open $B_r$ is open.
\frozen{KAM.actionNbhd}
\source{BM11 p.~4 ($D_s$); track K, \Kref{bp:p0:domains}.}
\end{definition}

\begin{definition}[Frequency map]\label{bpN:p0:freq}
For $h:\C^n\to\C$ and $I\in\R^n$, $\omega_h(I)=\nabla h(I)\in\R^n$ is the vector with
$\omega_h(I)_i=D\bigl(J\mapsto\operatorname{Re}h(J)\bigr)(I)[e_i]$, the real partial derivatives of
$\operatorname{Re}h$ on $\R^n$ at $I$. It is $0$ if $\operatorname{Re}h|_{\R^n}$ is not differentiable at $I$
(Lean totality, Remark~\ref{bpN:p0:totality}).
\frozen{KAM.frequencyMap}
\source{BM11 p.~1; Pos93 p.~190 ($\omega(I)=\partial_Ih(I)$).}
\end{definition}

\begin{definition}[Hamiltonian vector field]\label{bpN:p0:vf}
For $H:\R^n\times\R^n\to\R$ and $z=(\theta,I)$, $X_H(z)=\bigl((\partial_{I_i}H(z))_i,\,(-\partial_{\theta_i}H(z))_i\bigr)$,
where $\partial_{I_i}H(z)=DH(z)[(0,e_i)]$ and $\partial_{\theta_i}H(z)=DH(z)[(e_i,0)]$ are the Fr\'echet
derivatives in the Lean sense ($0$ if $H$ is not differentiable at $z$).
\frozen{KAM.hamiltonianVectorField}
\source{BM11 p.~1; track K, \Kref{bp:p0:hamiltonian}.}
\end{definition}

\begin{definition}[Real Hamiltonian]\label{bpN:p0:realH}
For $h:\C^n\to\C$, $f:\C^n\times\C^n\to\C$ and real $z=(\theta,I)$,
$H_{h,f}(z)=\operatorname{Re}\bigl(h(I)+f(\theta,I)\bigr)$, with $\theta,I$ viewed in $\C^n$.
\frozen{KAM.realHamiltonian}
\source{BM11 p.~5 ($H=h+f$).}
\end{definition}

\begin{remark}[Strip and neighbourhood notation]\label{bpN:p0:strip}
\Lean{KAM.complexStrip}\,$n\,s=T_s$ as above; for $n>0$ it is empty when $s\le0$, and for $n=0$ it is the
whole (one-point) space $\C^0$ for every $s$. No hypothesis $n>0$ is added to N-ZERO or N-QC.
\frozen{KAM.complexStrip}
\end{remark}

\subsection{Frozen track-N definitions}\label{bpN:p0:defs}

\begin{definition}[Action ball]\label{bpN:p0:ball}
$B(0,R)=\{I\in\R^n:\abs I<R\}$, the open ball of the max norm; written $\texttt{actionBall}\,n\,R$.
\frozen{KAM.Nekhoroshev.\allowbreak actionBall}
\source{BM11 p.~4.}
\end{definition}

\begin{definition}[Complex domain]\label{bpN:p0:cdomain}
$D_s(R)=T_s\times B(0,R)_s\subset\C^n\times\C^n$; \texttt{complexDomain}\,$n\,R\,s$. It is the
period-one lift of BM11's $D_s$, with a single width $s$ for angles and actions. It is open.
\frozen{KAM.Nekhoroshev.\allowbreak complexDomain}
\source{BM11 p.~4 ($D_s$); ours for the lift.}
\end{definition}

\begin{definition}[Real action domain]\label{bpN:p0:rdomain}
$\mathcal B_s(R)=\{I\in\R^n: I\in B(0,R)_s\}$; \texttt{realActionDomain}\,$n\,R\,s$.
\frozen{KAM.Nekhoroshev.\allowbreak realActionDomain}
\source{BM11 p.~4 ($B_s$).}
\end{definition}

\begin{remark}[Computing $\mathcal B_s(R)$]\label{bpN:p0:rdomain-eq}
Since $\abs{I-J}$ for real $I,J$ is the real max norm, $\mathcal B_s(R)=\{I\in\R^n:\exists J,\ \abs J<R,\
\abs{I-J}<s\}=B(0,R+s)$ for $R,s>0$. This identity is ours and is proved from the definitions (it is not stated in the
source); it shows that $B(0,R)\subset\mathcal B_s(R)$ and that $\mathcal B_s(R)$ is open.
\end{remark}

\begin{definition}[Real part of the integrable part]\label{bpN:p0:realpart}
For $h:\C^n\to\C$ and $I\in\R^n$, $h_{\R}(I)=\operatorname{Re}h(I)$; \texttt{realPart}\,$h$.
\frozen{KAM.Nekhoroshev.\allowbreak realPart}
\source{BM11 p.~4 ($h:B_s\to\R$).}
\end{definition}

\begin{definition}[Analytic integrable part, $h\in A_s(D)$]\label{bpN:p0:integrable}
$\texttt{IsAnalyticIntegrable}\,R\,s\,h$ holds if
\begin{enumerate}[label=(\roman*)]
\item $h$ is complex differentiable on $B(0,R)_s$;
\item $h$ is continuous on $\overline{B(0,R)_s}$ (closure in $\C^n$);
\item there is $A\ge0$ with $\abs{h(I)}\le A$ for all $I\in B(0,R)_s$;
\item $\operatorname{Im}h(I)=0$ for all $I\in\mathcal B_s(R)$;
\item $h_\R$ is of class $C^3$ on $\mathcal B_s(R)$ (real Fr\'echet sense; the set is open).
\end{enumerate}
No numerical bound on $h$ is fixed. Clause (v) is redundant by holomorphy (\Kref{bp:p1:holomorphic-regularity},
\Kref{bp:p1:real-restriction}) and records the regularity that guards every real derivative.
\frozen{KAM.Nekhoroshev.\allowbreak IsAnalyticIntegrable}
\source{BM11 p.~4 ($A_s(D)$, integrable case, with continuity on the closure).}
\end{definition}

\begin{definition}[Analytic perturbation of size $\varepsilon$]\label{bpN:p0:pert}
$\texttt{IsAnalyticPerturbation}\,R\,s\,\varepsilon\,f$ holds if $\varepsilon\ge0$ and
\begin{enumerate}[label=(\roman*)]
\item $f$ is complex differentiable on $D_s(R)$ and continuous on $\overline{D_s(R)}$;
\item $f(\theta+k,I)=f(\theta,I)$ for all $(\theta,I)\in D_s(R)$ and $k\in\Z^n$;
\item $\operatorname{Im}f(\theta,I)=0$ for all $\theta\in\R^n$ and $I\in\mathcal B_s(R)$;
\item $\abs{f(z)}\le\varepsilon$ for all $z\in D_s(R)$ (complex absolute value).
\end{enumerate}
The size is measured by $\le\varepsilon$, not the source's strict bound; the source is applied with $2\varepsilon$.
For $\varepsilon=0$ this forces $f=0$ on $D_s(R)$.
\frozen{KAM.Nekhoroshev.\allowbreak IsAnalyticPerturbation}
\source{BM11 pp.~4--5 ($f\in A_s(D)$, $\abs f_s<\varepsilon$); deviation: $\le\varepsilon$ (ruling Q2, D18).}
\end{definition}

\begin{definition}[Derivative bounds $B(M)$]\label{bpN:p0:bm}
For $U\subset\R^n$, $h:\C^n\to\C$ and $M\in\R$, $\texttt{HasDerivativeBounds}\,U\,h\,M$ holds if $M>0$, $U$ is open,
$h_\R\in C^3(U)$, and for all $I\in U$, all $k\in\{1,2,3\}$ and all $(i_1,\dots,i_k)\in\{1,\dots,n\}^k$,
\[\bigl\lvert D^kh_\R(I)[e_{i_1},\dots,e_{i_k}]\bigr\rvert\le M.\]
These are all ordered coordinate partials of orders one to three; for $C^3$ functions they coincide with the multi-index
partials $\partial^kh$ of the source (symmetry, \Kref{bp:p1:ext-schwarz}). The bound is on coordinate partials, not
on operator norms.
\frozen{KAM.Nekhoroshev.\allowbreak HasDerivativeBounds}
\source{BM11 p.~4, $B(M)$.}
\end{definition}

\begin{definition}[Strict quasi-convexity $QC(m)$]\label{bpN:p0:qc}
For $U\subset\R^n$, $h:\C^n\to\C$, $m\in\R$, $\texttt{IsStrictlyQuasiConvex}\,U\,h\,m$ holds if $m>0$, $U$ is open,
$h_\R\in C^2(U)$, and for every $I\in U$: $\omega_h(I)\ne0$ and
\[\sum_i\omega_h(I)_iv_i=0\ \Longrightarrow\ m\abs v^2\le D^2h_\R(I)[v,v]\qquad(v\in\R^n),\]
with $\abs v$ the max norm on the left and the Euclidean dot product in the hypothesis.
\frozen{KAM.Nekhoroshev.\allowbreak IsStrictlyQuasiConvex}
\source{BM11 p.~4, $QC(m)$ (the source's norm on the right is the max norm, orthogonality is Euclidean).}
\end{definition}

\begin{definition}[Complex Hessian bound]\label{bpN:p0:hess}
$\texttt{HasComplexHessianBound}\,R\,s\,h\,M_{\C}$ holds if $M_{\C}>0$, $h$ is $C^2$ in the complex sense on $B(0,R)_s$
and $\norm{D^2h(z)}\le M_{\C}$ for all $z\in B(0,R)_s$, the norm being that of $D(Dh)(z)\in\mathcal L(\C^n,\mathcal L(\C^n,\C))$
for the max norm. Every entry $\abs{\partial_i\partial_jh}$ is then $\le M_{\C}$, and the Euclidean operator norm of the
Hessian matrix is $\le nM_{\C}$. This is not the spectral radius of a complex symmetric matrix.
\frozen{KAM.Nekhoroshev.\allowbreak HasComplexHessianBound}
\source{Pos93 p.~190, ``The Hessian of the integrable hamiltonian $h$'' (conversion by $n$ recorded in
\texttt{\seqsplit{docs/tracks/N/sources.md}}, anchor POS190-HESS).}
\end{definition}

\begin{definition}[Solutions on a time set]\label{bpN:p0:sol}
For $H:\R^n\times\R^n\to\R$, $U\subset\R^n$, $J\subset\R$ and $z:\R\to\R^n\times\R^n$,
$\texttt{IsSolutionOn}\,H\,U\,J\,z$ holds if for every $t\in J$: $z(t)_2\in U$, $H$ is differentiable at $z(t)$, and
$z$ is differentiable at $t$ (two-sided, also at endpoints of $J$, since $z$ is defined on $\R$) with
$z'(t)=X_H(z(t))$.
\frozen{KAM.Nekhoroshev.\allowbreak IsSolutionOn}
\source{BM11 Theorem~2.1 p.~5 and the maximal-solution argument p.~13.}
\end{definition}

\begin{definition}[The integrable orbit]\label{bpN:p0:orbit}
For $h:\C^n\to\C$, $\theta_0,I_0\in\R^n$ and $t\in\R$:
$\gamma_{\theta_0,I_0}(t)=(\theta_0+t\,\omega_h(I_0),\,I_0)$.
\frozen{KAM.Nekhoroshev.\allowbreak integrableOrbit}
\source{BM11 p.~1 ($H(\theta,I)=h(I)$).}
\end{definition}

\begin{remark}[Lean totality]\label{bpN:p0:totality}
The frozen definitions and theorems use \Lean{fderiv} (which is $0$ at a point of non-differentiability),
\Lean{iteratedFDeriv}, the operator norm, \Lean{rpow} ($\varepsilon^{\pm1/(2n)}$), \Lean{Real.exp} and the max norm;
there is no matrix inverse, no Bochner integral and no division by a quantity that may vanish. The guards are:
\begin{enumerate}[label=(\alph*)]
\item In N-HEAD-C and N-ZERO, $\omega_h$ and $D^2h_\R$ are used only on the open set $\mathcal B_s(R)$ where $h_\R$
  is $C^3$ (Definitions~\ref{bpN:p0:integrable}(v), \ref{bpN:p0:bm}, \ref{bpN:p0:qc}), so the derivatives are the
  genuine ones. N-QC uses them on its arbitrary open $U$, where its hypothesis gives $C^2$ (item (e)). The
  definitions \texttt{frequencyMap}, \texttt{integrableOrbit} themselves are total and keep Lean's default
  values outside the differentiability locus.
\item $D^2h$ for complex $h$ is used only on $B(0,R)_s$, open, where $h$ is $C^2$ (Definition~\ref{bpN:p0:hess}).
\item $X_H$ is used only inside \texttt{IsSolutionOn}, which requires differentiability of $H$ at every visited point;
  in N-ZERO and N-HEAD-C this is a proof obligation: $H_{h,f}$ is real-differentiable on
  $\R^n\times\mathcal B_s(R)\supseteq\R^n\times B(0,R)$ (angles first, actions restricted) because $h,f$ are holomorphic on open sets containing the real
  points (\Kref{bp:p1:hamiltonian-real}).
\item $\varepsilon^{-1/(2n)}$ and $\varepsilon^{1/(2n)}$ occur only with $\varepsilon>0$ in N-HEAD-C, and $2n\ne0$.
\item In N-QC the map $\Psi(I',\lambda)=(h_\R(I'),\lambda\omega_h(I'))$ is differentiable at $(I,\lambda)$ because
  $h_\R\in C^2$ on the open set $U\ni I$; \Lean{fderiv} is then the true derivative.
\end{enumerate}
\end{remark}

\subsection{The frozen theorems}\label{bpN:p0:theorems}

\begin{theorem}[N-HEAD-C: exponential stability, complex Hessian bound]\label{bpN:p0:head}
Let $n\ge3$ and $R,s,M,m,M_{\C},\omega_{\min}>0$. Then there are $C,T,c,\varepsilon_0>0$ with
\[C\,\varepsilon_0^{1/(2n)}\le R/4\]
such that the following holds for every $h:\C^n\to\C$, $f:\C^n\times\C^n\to\C$ and $\varepsilon\in\R$ with
\begin{enumerate}[label=(H\arabic*)]
\item $\texttt{IsAnalyticIntegrable}\,R\,s\,h$ (Definition~\ref{bpN:p0:integrable});
\item $\texttt{HasComplexHessianBound}\,R\,s\,h\,M_{\C}$ (Definition~\ref{bpN:p0:hess});
\item $\omega_{\min}\le\abs{\omega_h(I)}$ for every $I\in B(0,R)$ (max norm);
\item $\texttt{IsStrictlyQuasiConvex}\,\mathcal B_s(R)\,h\,m$ (Definition~\ref{bpN:p0:qc});
\item $\texttt{HasDerivativeBounds}\,\mathcal B_s(R)\,h\,M$ (Definition~\ref{bpN:p0:bm});
\item $\texttt{IsAnalyticPerturbation}\,R\,s\,\varepsilon\,f$ (Definition~\ref{bpN:p0:pert});
\item $0<\varepsilon\le\varepsilon_0$.
\end{enumerate}
For every $\theta_0\in\R^n$ and $I_0\in B(0,R/2)$, with $\tau:=T\exp\bigl(c\,\varepsilon^{-1/(2n)}\bigr)$ and
$H=H_{h,f}$, there is $z:\R\to\R^n\times\R^n$ with
\begin{enumerate}[label=(\alph*)]
\item $z(0)=(\theta_0,I_0)$;
\item $\texttt{IsSolutionOn}\,H\,B(0,R)\,[-\tau,\tau]\,z$;
\item $\abs{z(t)_2-I_0}\le C\varepsilon^{1/(2n)}$ for all $t$ with $\abs t\le\tau$;
\item (uniqueness) for every $L\in[0,\tau]$ and every $y:\R\to\R^n\times\R^n$ with $y(0)=(\theta_0,I_0)$ and
  $\texttt{IsSolutionOn}\,H\,B(0,R)\,[-L,L]\,y$: $y(t)=z(t)$ for all $t\in[-L,L]$.
\end{enumerate}
The order of choice is: $n$; $R,s,M,m,M_{\C},\omega_{\min}$; then $C,T,c,\varepsilon_0$, depending only on these
seven data; then $h,f,\varepsilon$, then $\theta_0,I_0$; the solution $z$ may depend on all of them. There is no
dependence on $A$ in Definition~\ref{bpN:p0:integrable}(iii). In (H4)--(H5) the set is $\mathcal B_s(R)=B(0,R+s)$
(Remark~\ref{bpN:p0:rdomain-eq}), in (H3) the smaller ball $B(0,R)$.
\frozen{KAM.Nekhoroshev.\allowbreak quasiConvex\_\allowbreak nekhoroshev\_\allowbreak complexBound}
\source{Pos93 Theorem~1, pp.~190--191 (with the conversions and anchors of \texttt{\seqsplit{docs/tracks/N/sources.md}}); exponent $1/(2n)$
as BM11 p.~5 with $\delta=(2n(n-1))^{-1}$; ruling D18 Q2, Q4; deviation D20: the extra inputs $M_{\C}$ and $\omega_{\min}$.}
\sourcegap{The printed theorem is for a $2\pi$-periodic torus, Euclidean norms and exceptional initial frequencies
$\norm{\omega(I_0)}\le mr_0/8$; the conversion to the frozen statement (rescaling $H\mapsto2\pi H(\varphi/2\pi,I)$,
choice of $r_0,s_0$ from the seven data, exclusion of the exceptional branch by $\omega_{\min}$, existence and
uniqueness of the solution) is given only in \texttt{sources.md} and is proved in Part~10 of \texttt{PLAN.md}.}
\end{theorem}

\begin{theorem}[N-ZERO: exact integrable dynamics]\label{bpN:p0:zero}
Let $n\in\N$, $R,s>0$, $h$ with $\texttt{IsAnalyticIntegrable}\,R\,s\,h$ and $f$ with
$\texttt{IsAnalyticPerturbation}\,R\,s\,0\,f$. For every $\theta_0\in\R^n$ and $I_0\in B(0,R/2)$, with
$\gamma=\gamma_{\theta_0,I_0}$ (Definition~\ref{bpN:p0:orbit}):
\[\gamma(0)=(\theta_0,I_0),\qquad \texttt{IsSolutionOn}\,H_{h,f}\,B(0,R)\,\R\,\gamma,\qquad \gamma(t)_2=I_0\ \ (t\in\R).\]
\frozen{KAM.Nekhoroshev.\allowbreak zero\_perturbation}
\source{Ours (supplement to BM11 Theorem~2.1, p.~5, whose strict bound excludes $\varepsilon=0$); from $H=h(I)$, BM11 p.~1.}
\end{theorem}

\begin{theorem}[N-QC: nonsingular isoenergetic derivative]\label{bpN:p0:qcthm}
Let $n\in\N$, $U\subset\R^n$, $h:\C^n\to\C$, $m\in\R$ with $\texttt{IsStrictlyQuasiConvex}\,U\,h\,m$. For $I\in U$ and
$\lambda>0$, the derivative at $(I,\lambda)$ of
\[\Psi:\R^n\times\R\to\R\times\R^n,\qquad\Psi(I',\lambda')=\bigl(h_\R(I'),\ \lambda'\omega_h(I')\bigr)\]
is a bijective linear map $\R^n\times\R\to\R\times\R^n$. Explicitly it is
$D\Psi(I,\lambda)[v,u]=\bigl(\omega_h(I)\cdot v,\ u\,\omega_h(I)+\lambda D^2h_\R(I)v\bigr)$.
No uniform radius of invertibility is asserted.
\frozen{KAM.Nekhoroshev.\allowbreak quasiConvex\_\allowbreak isoenergetic\_\allowbreak derivative}
\source{BM11 Lemma~3.3 and its proof, pp.~9--10 (only the derivative assertion).}
\end{theorem}

\subsection{D25 additions: the partial (slow--fast) convex theorem}\label{bpN:p0:partial-sec}

Owner approval D25 (27 October 2026) froze one further theorem, N-PARTIAL-CONVEX, and seven definitions. The Lean
declarations are authoritative (\texttt{\seqsplit{KAM/Frozen/Nekhoroshev/Theorems/PartialConvexNekhoroshev.lean}},
\texttt{\seqsplit{KAM/Frozen/Nekhoroshev/Defs/\{ComplexifyPhase,InnerCore,PartialComplexDomain,HasDefiniteFastPart,IsPartialPerturbation,PartialTimeExponent,PartialDriftExponent\}.lean}}).
Notation: $\mathrm{Ph}_d=\R^d\times\R^d$ (Lean \Lean{KAM.Planetary.Phase}\,$d$, positions first, then momenta) with the product of max norms,
$\norm{(x,y)}=\max(\abs x,\abs y)$. For $v\in\R^{n+k}$ (or $\R^{n+k}$-valued vectors), $\mathrm{fast}(v)=(v_1,\dots,v_n)$ is the first block and
$\mathrm{slow}(v)=(v_{n+1},\dots,v_{n+k})$ the second block (\Lean{KAM.Planetary.fast}, \Lean{KAM.Planetary.slow}).
A phase point of the partial system is $z=(\varphi,\mathsf I)\in\mathrm{Ph}_{n+k}$: $\varphi=(\varphi_{\rm f},q)$ with $n$ fast angles $\varphi_{\rm f}=\mathrm{fast}(\varphi)$ and
$k$ slow positions $q=\mathrm{slow}(\varphi)$; $\mathsf I=(I,p)$ with $n$ fast actions $I=\mathrm{fast}(\mathsf I)$ and $k$ slow momenta $p=\mathrm{slow}(\mathsf I)$. The slow pair is $(q,p)=\bigl(\mathrm{slow}(\varphi),\mathrm{slow}(\mathsf I)\bigr)\in\mathrm{Ph}_k$.
The fast angles have period $2\pi$ (\Lean{KAM.Planetary.fastShift}\,$\ell=(2\pi\ell,0)$ in the fast block).

\begin{definition}[\Lean{complexifyPhase}]\label{bpN:p0:partial-complexify}
For $z=(x,y)\in\mathrm{Ph}_d$: $\mathrm{cx}(z)=\bigl((x_i)_i,(y_i)_i\bigr)\in\C^d\times\C^d$, componentwise inclusion $\R\subset\C$; no choices.
\frozen{KAM.Nekhoroshev.\allowbreak complexifyPhase}
\source{Ours (notation for Nek77 \S4.3, p.~29: ``for real values of the arguments $H$ is real-valued'').}
\end{definition}

\begin{definition}[\Lean{innerCore}: closed max-norm erosion]\label{bpN:p0:partial-core}
For a normed group $E$, $U\subset E$ and $d\in\R$: $\operatorname{core}(U,d)=\{x\in U:\ \bar B(x,d)\subset U\}$ (closed ball of the norm of $E$). Used only with $d>0$ (for $d<0$ the ball is empty and $\operatorname{core}(U,d)=U$).
It is used on $E=\R^n$ (max norm) and $E=\mathrm{Ph}_k$ (product of max norms).
\frozen{KAM.Nekhoroshev.\allowbreak innerCore}
\source{Nek77 \S4.2, p.~29 ($X-\varepsilon$: points of $X$ whose $\varepsilon$-neighbourhoods lie in $X$); here closed balls of the max norm, a stronger margin than the Euclidean erosion of the source.}
\end{definition}

\begin{definition}[\Lean{partialComplexDomain}]\label{bpN:p0:partial-domain}
For $G\subset\R^n$, $D\subset\mathrm{Ph}_k$ and $\rho\in\R$,
\[\mathcal F(G,D,\rho)=\Bigl\{(\varphi,\mathsf I)\in\C^{n+k}\times\C^{n+k}:\ \mathrm{fast}(\operatorname{Re}\mathsf I)\in G,\ \bigl(\mathrm{slow}(\operatorname{Re}\varphi),\mathrm{slow}(\operatorname{Re}\mathsf I)\bigr)\in D,\ \varphi,\mathsf I\in T^{n+k}_\rho\Bigr\},\]
$T^{N}_\rho=\{u\in\C^N:\abs{\operatorname{Im}u_i}<\rho\ \forall i\}$ (Lean \Lean{KAM.complexStrip}). The fast real angles and the slow real positions carry no condition other than $(\mathrm{slow}\operatorname{Re}\varphi,\mathrm{slow}\operatorname{Re}\mathsf I)\in D$.
\frozen{KAM.Nekhoroshev.\allowbreak partialComplexDomain}
\source{Nek77 \S4.3.A, p.~29 (complex tube over arbitrary fast/slow real domains $G$, $D$; there the imaginary bounds are closed, $\le\rho$, and the slow variables are $(p,q)\in D\subset\R^{2(n-s)}$).}
\end{definition}

\begin{definition}[\Lean{HasDefiniteFastPart}]\label{bpN:p0:partial-definite}
For $G\subset\R^n$, $\rho,g,m\in\R$ and $h:\C^n\to\C$, with $U=\{I\in\R^n:\ I\in G_\rho\}$ ($G_\rho=\texttt{actionNbhd}\,G\,\rho$, Definition~\ref{bpN:p0:actionnbhd}; its real points):
$\texttt{HasDefiniteFastPart}\,G\,\rho\,g\,m\,h$\allowbreak\ holds iff $g>0$, $m>0$, the function $h_\R=\operatorname{Re}h|_{\R^n}$ is $C^3$ on $U$ (Lean \Lean{ContDiffOn}\,$\R\,3$, with $U$ the set above) and there is $\sigma\in\{1,-1\}$ with, for every $I\in U$:
\[\abs{\omega_h(I)}\ge g\qquad\text{and}\qquad \forall v\in\R^n:\quad m\sum_iv_i^2\le\sigma\,D^2h_\R(I)[v,v],\]
where $\omega_h=\nabla h_\R$ (Definition~\ref{bpN:p0:freq}; the norm of $\omega_h(I)\in\R^n$ is the max norm), and $D^2h_\R(I)[v,v]$ is the second \Lean{fderiv} of $h_\R$ at $I$ applied to $v$ and $v$ (Lean \Lean{fderiv}\,$\R$ twice).
\frozen{KAM.Nekhoroshev.\allowbreak HasDefiniteFastPart}
\source{Specialisation of Nek77 \S1.18 and \S4.1 (index-one steep class) to a definite fast Hessian of one global sign $\sigma$; deviation D-N-ADD-1 (\texttt{docs/tracks/N/REVIEW-additions.md}); not a source definition.}
\end{definition}

\begin{definition}[\Lean{IsPartialPerturbation}]\label{bpN:p0:partial-pert}
For $G\subset\R^n$, $D\subset\mathrm{Ph}_k$, $\rho\in\R$ and $P:\C^{n+k}\times\C^{n+k}\to\C$, $\texttt{IsPartialPerturbation}\,G\,D\,\rho\,P$ holds iff
\begin{enumerate}[label=(P\arabic*)]
\item $P$ is complex differentiable on $\mathcal F(G,D,\rho)$ (Lean \Lean{DifferentiableOn}\,$\C$);
\item $\abs{P(z)}\le1$ for $z\in\mathcal F(G,D,\rho)$;
\item for real $z\in\mathrm{Ph}_{n+k}$ with $\mathrm{cx}(z)\in\mathcal F(G,D,\rho)$: $\operatorname{Im}P(\mathrm{cx}(z))=0$;
\item for every $z\in\mathcal F(G,D,\rho)$ and $\ell\in\Z^n$: $P\bigl(z+\mathrm{cx}(\mathrm{fastShift}\,\ell)\bigr)=P(z)$, i.e.\ $P$ is $2\pi$-periodic in each fast angle; there is no periodicity in the slow variables.
\end{enumerate}
\frozen{KAM.Nekhoroshev.\allowbreak IsPartialPerturbation}
\source{Nek77 \S4.3 and Remark 4.5.A, p.~30 (the perturbation $H_1$ is analytic on $F$, $2\pi$-periodic in $\varphi$, real for real arguments; its size $M=\sup\abs{H_1}$ is $\le\mu$ for $H_1=\mu P$).}
\end{definition}

\begin{definition}[\Lean{partialTimeExponent}, \Lean{partialDriftExponent}]\label{bpN:p0:partial-exponents}
For $n\in\N$: $a_n=\dfrac2{6n^2-3n+14}$ and $b_n=\dfrac3{6n^2-3n+14}=\tfrac32a_n$ (real arithmetic; for $n\ge2$ the denominator is $\ge 6\cdot4-6+14=32>0$).
\frozen{KAM.Nekhoroshev.\allowbreak partialTimeExponent, KAM.Nekhoroshev.\allowbreak partialDriftExponent}
\source{Nek77 (4.1), (4.4), (4.5), p.~29--30, with all steepness indices $\alpha_i=1$ and $s=n$ (the formulas (4.4), (4.5) are images that the text layer of the scan loses; the frozen definitions and \texttt{docs/tracks/N/statement.md} are used); for $n=2$: $a=\frac1{16}$, $b=\frac3{32}$ (\texttt{docs/tracks/N/statement.md}).}
\end{definition}

\begin{theorem}[N-PARTIAL-CONVEX: exponential stability of the fast actions, definite fast part, slow--fast]\label{bpN:p0:partial}
Let $n\ge2$, $k\ge1$, $G,G'\subset\R^n$ open and convex, $D\subset\mathrm{Ph}_k$ open, $\rho,g,m,M_{\C}\in\R$ with $\rho>0$, $M_{\C}>0$ and
\[\bar B(I,\rho)\subset G'\quad\text{for every }I\in G\qquad(\text{closed max-norm ball}),\]
and let $h:\C^n\to\C$ with: $h$ complex differentiable and complex $C^2$ on $G'_\rho$; $\operatorname{Im}h(I)=0$ for every real $I$ with $I\in G'_\rho$; $\norm{D^2h(z)}\le M_{\C}$ for every $z\in G'_\rho$ (operator norm of the complex Hessian, max norms); and $\texttt{HasDefiniteFastPart}\,G'\,\rho\,g\,m\,h$. Then there is $\mu_0>0$ such that for all $\mu\in(0,\mu_0)$, all $P$ with $\texttt{IsPartialPerturbation}\,G'\,D\,\rho\,P$, all $C>0$ and all $z:\R\to\mathrm{Ph}_{n+k}$ with
\begin{enumerate}[label=(Z\arabic*)]
\item $\texttt{IsSolutionOn}\;H_\mu\;\{\mathsf I:\mathrm{fast}(\mathsf I)\in G\}\;[0,C]\;z$, where $H_\mu(\varphi,\mathsf I)=\operatorname{Re}h\bigl(\mathrm{fast}(\mathsf I)\bigr)+\mu\operatorname{Re}P\bigl(\mathrm{cx}(\varphi,\mathsf I)\bigr)$ (Lean \Lean{KAM.Planetary.degenerateHamiltonian}\,$(\texttt{realPart }h)\,(\lambda\_\,w.\ \operatorname{Re}P(\mathrm{cx}\,w))\,\mu$, Definition~\ref{bpN:p0:realpart});
\item $\mathrm{fast}(z(0)_2)\in\operatorname{core}(G,\mu^{b_n})$,
\end{enumerate}
and every $t\in\bigl[0,\min\bigl(C,\ \mu^{-1}\exp(\mu^{-a_n})\bigr)\bigr]$ with
\[\Bigl(\mathrm{slow}(z(u)_1),\mathrm{slow}(z(u)_2)\Bigr)\in\operatorname{core}\bigl(D,\mu^{b_n}\bigr)\quad\text{for all }u\in[0,t],\]
one has
\[\bigl|\mathrm{fast}(z(t)_2)-\mathrm{fast}(z(0)_2)\bigr|\le\tfrac12\mu^{b_n}.\]
The order of choice is: $n,k$; $G,G',D,\rho,g,m,M_{\C},h$ (together with their hypotheses); then $\mu_0$; then $\mu,P,C,z$; then $t$. The slow confinement is tested only on $[0,t]$ (never assumed up to the exponential time).
\frozen{KAM.Nekhoroshev.\allowbreak partialConvex\_\allowbreak nekhoroshev}
\source{Nek77 Theorem 4.4, p.~30, with Remarks 4.5.A--B; proof specialised to a definite fast Hessian as in \texttt{docs/tracks/N/notes.md}, F-NADD-1 (Lemma 10.6, p.~47, replaced by its definite-Hessian endpoint argument); deviation D-N-ADD-1 (collar $G\Subset G'$, $2\pi$ period, stopped quantifier, maximum-norm erosion).}
\sourcegap{The printed theorem assumes steepness of $H_0$ on the whole source domain with explicit coefficients, a Euclidean erosion, closed imaginary bounds, and the confinement of $(p,q)$ on the whole $[0,C]$; the frozen statement (collar, closed max-norm erosion, confinement tested on $[0,t]$ only, imaginary width $\rho$ open) is proved here via the route of Part~\ref{bpN:p11}. The initial erosion $\operatorname{core}(G,\mu^{b_n})$ is used only through $\mathrm{fast}(z(0)_2)\in G$.}
\end{theorem}

\begin{remark}[Lean totality for the D25 objects]\label{bpN:p0:partial-totality}
(a) $\mu^{b_n}$, $\mu^{-a_n}$, $\mu^{-1}$ occur only for $\mu>0$ (rpow of a positive base), and $6n^2-3n+14>0$. (b) $\omega_h$, $D^2h_\R$ in \texttt{HasDefiniteFastPart} are guarded: $h_\R$ is $C^3$ on the real slice $U$ of the open set $G'_\rho$, which is open in $\R^n$ ($G'$ open), so the \Lean{fderiv} are true derivatives at every $I\in U$ (Lemma~\ref{bpN:p1:holo-real} for the relation with the complex derivatives). (c) $\texttt{IsSolutionOn}$ asserts differentiability of $H_\mu$ at $z(u)$, $u\in[0,C]$; where $\mathrm{slow}(z(u))\notin D$ the value $P(\mathrm{cx}\,w)$ is the arbitrary value of the function $P$ outside $\mathcal F$, and the theorem uses the solution only for $u\in[0,t]$ where the slow variables lie in $\operatorname{core}(D,\mu^{b_n})\subset D$, an open set on which $H_\mu$ is real analytic. (d) If $\operatorname{core}(D,\mu^{b_n})$ is empty the confinement hypothesis fails for $t\ge0$ and the statement is vacuous. (e) $\texttt{fderiv}\,\C\,(\texttt{fderiv}\,\C\,h)$ is the true second derivative on the open set $G'_\rho$ where $h$ is complex $C^2$.
\end{remark}

\subsection{D28 additions: multi-frequency averaging}\label{bpN:p0:avg-sec}

Owner approval D28 (27 October 2026) froze two further theorems, N-AVERAGING-MEAN and N-AVERAGING-MOST, and two definitions
(\texttt{\seqsplit{KAM/Frozen/Nekhoroshev/Theorems/\{AveragingMean,AveragingMost\}.lean}},
\texttt{\seqsplit{KAM/Frozen/Nekhoroshev/Defs/\{HasAveragingFrequencyRank,AveragingError\}.lean}}).
The Lean text is authoritative. Notation of the Lean statements: a \emph{phase point} is $x=(\theta,I)\in\R^n\times\R^n$, the action is
\texttt{x.2}$=x_2=I$, norms are max norms, $\texttt{vol}$ is Lebesgue measure on $\R^n\times\R^n$ (``unnormalised phase volume'': no
division by $\operatorname{vol}K$ and no probability), and for $\varepsilon>0$ the Hamiltonian is
$H_\varepsilon=H_{h,\varepsilon f}=\operatorname{Re}\bigl(h(I)+\varepsilon f(\theta,I)\bigr)$ (Definition~\ref{bpN:p0:realH} with the perturbation
$w\mapsto(\varepsilon:\C)\,f(w)$). With the averaged action orbit $J(t)\equiv I_0$ (Remark~\ref{bpN:p12:averaged}) the error
of the averaging method is $\sup_{0\le t\le1/\varepsilon}\abs{I(t)-I_0}$.

\begin{definition}[\Lean{HasAveragingFrequencyRank}: the two rank alternatives]\label{bpN:p0:avg-rank}
For $U\subset\R^n$ and $h:\C^n\to\C$ write $\omega=\omega_h$ (Definition~\ref{bpN:p0:freq}). $\texttt{HasAveragingFrequencyRank}\,U\,h$ holds iff
\[\text{(A)}\ \ \forall I\in U:\ D\omega(I):\R^n\to\R^n\ \text{is surjective}\qquad\text{or}\qquad
\text{(B)}\ \ \forall I\in U:\ \omega(I)\ne0\ \text{and}\ (v,c)\mapsto D\omega(I)v+c\,\omega(I),\ \R^n\times\R\to\R^n,\ \text{is surjective,}\]
with $D\omega(I)=\texttt{fderiv}\,\R\,\omega\,I$. The alternative is global: the same branch (A) or (B) holds at every point of $U$
(the quantifier $\forall I\in U$ is inside each branch). (B) says that the projective class $(\omega_1:\dots:\omega_n)$ has differential of rank $n-1$ at $I$
(its image is $(\operatorname{range}D\omega(I)+\R\omega(I))/\R\omega(I)$); no chart $\omega_1\ne0$ is chosen. Branch (A) may hold at a point where $\omega=0$.
\frozen{KAM.Nekhoroshev.\allowbreak HasAveragingFrequencyRank}
\source{AKN06 Theorem~6.11, p.~243 (PDF p.~254): ``either the rank of the map $I\mapsto\omega(I)$ is equal to the number of frequencies, or the rank of the map $I\mapsto(\omega_1(I):\dots:\omega_m(I))$ is by one less than the number of frequencies''; Nei76 p.~1296 hypothesis~3$^\circ$ is the chart form.}
\end{definition}

\begin{definition}[\Lean{averagingError}]\label{bpN:p0:avg-error}
For $n\in\N$, $\varepsilon\in\R$, $\Phi:(\R^n\times\R^n)\to\R\to(\R^n\times\R^n)$ and $x=(\theta_0,I_0)$:
\[e_\varepsilon^\Phi(x)=\texttt{averagingError}\,\varepsilon\,\Phi\,x=\sup\Bigl\{\abs{(\Phi\,x\,t)_2-I_0}:\ t\in[0,1/\varepsilon]\Bigr\}\in\R,\]
the Lean real supremum \Lean{sSup} (which is $0$ for an empty or unbounded-above set).
\frozen{KAM.Nekhoroshev.\allowbreak averagingError}
\source{AKN06 (6.27), p.~243: $\max_{0\le t\le1/\varepsilon}\abs{I(t)-J(t)}$ with $J(t)=I_0$; Nei76 p.~1295 (4).}
\end{definition}

\begin{remark}[Lean totality for $\texttt{averagingError}$]\label{bpN:p0:avg-totality}
\Lean{sSup} on $\R$ is the true supremum only for a nonempty set bounded above; otherwise it is the junk value $0$. The theorems use the definition
only under the following guards, each verified in Part~\ref{bpN:p12}:
\begin{enumerate}[label=(\alph*)]
\item $\varepsilon>0$ (the quantifier $0<\varepsilon\le\varepsilon_0$), so $1/\varepsilon>0$ and $0\in[0,1/\varepsilon]$: the interval is a nonempty compact interval (for $\varepsilon<0$ it is empty, for $\varepsilon=0$ it is $\{0\}$ since $1/0=0$ in Lean);
\item the family is a genuine solution on $[0,1/\varepsilon]$, so $t\mapsto(\Phi x\,t)_2$ is continuous on the compact interval (a solution is differentiable at every $t\in J$), whence the image of $t\mapsto\abs{(\Phi x\,t)_2-x_2}$ is compact, nonempty and bounded; moreover $(\Phi x\,t)_2\in B(0,R)$ and $x_2\in B(0,R/2)$ bound the set by $3R/2$ (the bound is also available without continuity);
\item then the supremum is attained, is $\ge0$ (it is $\ge$ the value $0$ at $t=0$ when $\Phi x\,0=x$), and $\abs{(\Phi x\,t)_2-x_2}\le e^\Phi_\varepsilon(x)$ for every $t\in[0,1/\varepsilon]$.
\end{enumerate}
No division by a quantity that may vanish occurs: $1/\varepsilon$ with $\varepsilon>0$, and $\sqrt\varepsilon$ with $\varepsilon>0$, $A\sqrt\varepsilon/\kappa$ with $\kappa>0$, $A\sqrt\varepsilon/\eta$ with $\eta>0$. The derivatives behind $\omega_h$ and $D\omega_h$ in
Definition~\ref{bpN:p0:avg-rank} are genuine on $B(0,R)\subset\mathcal B_s(R)$ because $h_\R\in C^3$ there (Definition~\ref{bpN:p0:integrable}(v)), so $\omega_h$ is $C^2$ and \Lean{fderiv} is the true derivative. As a stand-alone predicate on an arbitrary $h$ or on $U=\emptyset$ the definition has its usual vacuity/off-locus behaviour; no theorem uses it there.
\end{remark}

\begin{theorem}[N-AVERAGING-MEAN: mean error and Markov form]\label{bpN:p0:avg-mean}
Let $n\ge2$, $R,s>0$, $h:\C^n\to\C$, $f:\C^n\times\C^n\to\C$ with $\texttt{IsAnalyticIntegrable}\,R\,s\,h$, $\texttt{IsAnalyticPerturbation}\,R\,s\,1\,f$,
$\texttt{HasAveragingFrequencyRank}\,B(0,R)\,h$, and $K\subset\R^n\times\R^n$ compact with $x_2\in B(0,R/2)$ for every $x\in K$. Then there are $A,\varepsilon_0>0$ such that for every $\varepsilon\in(0,\varepsilon_0]$ and
every $\Phi:(\R^n\times\R^n)\to\R\to(\R^n\times\R^n)$ with
\[\forall x\in K:\quad\Phi\,x\,0=x\quad\text{and}\quad\texttt{IsSolutionOn}\ H_{h,\varepsilon f}\ B(0,R)\ [0,1/\varepsilon]\ (\Phi\,x),\]
the function $e^\Phi_\varepsilon$ is integrable on $K$,
\[\int_K^{-}\operatorname{ofReal}\bigl(e^\Phi_\varepsilon(x)\bigr)\,dx\le\operatorname{ofReal}(A\sqrt\varepsilon)\qquad(\text{Lean \Lean{lintegral} over }K),\]
and for every $\kappa>0$ there is a measurable $E\subset K$ with $\operatorname{vol}E\le\kappa$ and $\abs{(\Phi\,x\,t)_2-x_2}\le A\sqrt\varepsilon/\kappa$ for all $x\in K\setminus E$ and $t\in[0,1/\varepsilon]$.
The order of choice: $n,R,s,h,f,K$ and their hypotheses; then $A,\varepsilon_0$; then $\varepsilon$, $\Phi$; then $\kappa$ and $E$. The flow premise is only existence and typing on $[0,1/\varepsilon]$; it contains no error bound and no exceptional set.
\frozen{KAM.Nekhoroshev.\allowbreak averaging\_\allowbreak mean}
\source{AKN06 Theorem~6.11 and Corollary~6.1 with (6.27), (6.28), p.~243 (printed pages; PDF $=$ printed $+11$); Nei76 pp.~1295--1298, Theorem~1; analytic specialisation, period-one angles, max norms and unnormalised volume as in \texttt{docs/tracks/N/statement.md}, \S N-AVERAGING v0.5.}
\sourcegap{Nei76 p.~1298 (17) writes $\sup_tz(t)=\int_{U'}h$, interchanging a supremum in time with an integral over initial data (ledger E-NAVG-NEI-17). The proof in Part~\ref{bpN:p12} bounds the integral of the maximum directly (Remark~\ref{bpN:p12:errata}).}
\end{theorem}

\begin{theorem}[N-AVERAGING-MOST: existence and small drift for most initial data]\label{bpN:p0:avg-most}
With the same $n,R,s,h,f,K$ and hypotheses as in Theorem~\ref{bpN:p0:avg-mean}, there are $A,\varepsilon_0,\eta_0>0$ with $\eta_0\le R/4$ such that for every $\varepsilon\in(0,\varepsilon_0]$ and every $\eta\in(0,\eta_0]$ there is a measurable $E\subset K$ with
\[\operatorname{vol}E\le\operatorname{ofReal}\bigl(A\sqrt\varepsilon/\eta\bigr),\]
and for every $x\in K\setminus E$ there is $z:\R\to\R^n\times\R^n$ with $z(0)=x$, $\texttt{IsSolutionOn}\ H_{h,\varepsilon f}\ B(0,R)\ [0,1/\varepsilon]\ z$ and $\abs{z(t)_2-x_2}\le\eta$ for all $t\in[0,1/\varepsilon]$.
The order of choice: the fixed data; then $A,\varepsilon_0,\eta_0$; then $\varepsilon$, then $\eta$, then $E$, then $x$ and $z$. Existence of the solution on the whole window and confinement are \emph{conclusions} on the good data, not premises.
\frozen{KAM.Nekhoroshev.\allowbreak averaging\_\allowbreak most}
\source{AKN06 p.~243 (6.28) and Nei76 Theorems~1--2, pp.~1295--1298 (the set $\Phi(\rho,\varepsilon)$ of p.~1296 includes existence on $[0,\varepsilon^{-1}]$ and its complement includes data that reach $\partial G$); AKN06 p.~3 footnote~1 (smoothness convention) and Remark~6.4, pp.~217--218 (averaged orbit away from the boundary).}
\sourcegap{As for Theorem~\ref{bpN:p0:avg-mean}; in addition the printed Theorem~2 of Nei76 imposes $\rho<\rho_1/2$ with an auxiliary collar $\rho_1$; the route of Part~\ref{bpN:p12} uses the exit radius $\rho=R/8$ and $\eta_0=R/8$ (a proof choice inside the existential constants; no strict inequality $\eta<\rho/2$ is needed there).}
\end{theorem}
